\documentclass[runningheads]{llncs}

\usepackage{graphicx}
\usepackage{booktabs}
\usepackage{hyperref}

\begin{document}

\title{The Brainrot Tracker: An Early Warning System for Doomscroll Susceptibility from Smartphone Kinematics and Interaction Behaviour}

\author{Author One \and Author Two \and Author Three}

\authorrunning{Author One et al.}

\institute{Vrije Universiteit Amsterdam, Amsterdam, The Netherlands}

\maketitle

\begin{abstract}
Compulsive engagement with algorithmic short-form video applications is a pervasive
behavioural phenomenon linked to cognitive fatigue, stress, and boredom. This paper
presents the \emph{Brainrot Tracker}, a dual-modal Early Warning System (EWS) that
predicts whether a user is about to initiate a doomscrolling session in the next
three-minute window. We continuously collect high-frequency smartphone kinematic data
(accelerometer and gyroscope at 10\,Hz) together with event-driven system-interaction logs
(target and auxiliary application launches) from three participants over the full duration
of the study. Each three-minute block is aggregated into engineered features, reframing the
extreme class-imbalance problem of precise event prediction as a tractable binary
classification task. We compare a classical gradient-boosted model against a temporal deep
learning model, and---crucially---assess generalisation \emph{across people} by training on
a subset of participants and testing on a held-out individual. <describe headline results>

\keywords{Quantified Self \and Behavioural Prediction \and Early Warning Systems \and
Time-Series Classification \and Deep Learning}
\end{abstract}


% -----------------------------------------------------------------------
\section{Research Question and Data}\label{sec:intro}
% -----------------------------------------------------------------------

\subsection{Problem Statement}

Short form video content is increasingly prevelant in the lives of the younger generation. It is rarely a pre-planned activity like other hobbies, but rather "doomscrolling" sessions arise seemingly out of coincidental circumstances. Our tracker aims to leverage data to predict when a user is about to open a short form video content app. While predicting the \emph{exact} second a user opens one of these applications is unfeasable as a machine learning problem, with the positive class (app opened)
representing fewer than 0.1\% of all observations in a continuous monitoring dataset -
producing severe class imbalance, we can reframe the problem as an \textbf{Early Warning System (EWS)}. Rather than predicting
the precise moment of failure, the model predicts whether a user is in a
\emph{Susceptibility Window} - a three-minute block immediately preceding a doomscrolling
session. This formulation is both more tractable
and more practically useful: a short warning horizon is still long enough for a behavioural
intervention (e.g.\ a mindfulness prompt or app lock) to be effective - ideally increasing the users awareness of the behavioural patterns that lead to this destructive behaviour.

\subsection{Research Questions}

This project is structured around two core research questions:

\begin{enumerate}
    \item \textbf{Predictability.} Can we predict brainrot time intervals? That is, using
          smartphone kinematic data (stillness, movement intensity) and system-interaction
          history aggregated over the current window, can we predict whether a user will
          enter a doomscrolling session in the subsequent window?
    \item \textbf{Transferability.} Does a learned model transfer across people? If a model
          is trained on some users' data, does its predictive performance hold when applied
          to a different user who was never seen during training?
\end{enumerate}

The second question directly satisfies the assignment's requirement that generalisation be
measured on \emph{another person}, rather than by randomly sampling time points from a
single dataset.

\subsection{Prediction Target}

The binary classification target is defined by a rolling three-minute look-ahead window
applied to the continuous time-series. Features computed during the current interval ($T$)
are mapped to a label derived from the subsequent interval ($T{+}1$). A block is assigned
\textbf{Class 1} (Susceptible State) if a verified TikTok or Instagram app-open event occurs
during the next three-minute interval, and \textbf{Class 0} (Baseline State) otherwise. This
window size was selected to balance label prevalence against the responsiveness of the
warning.

\subsection{Data Sources and Measurements}

Data are collected across two complementary modalities using the mobile devices of all
three authors. Each author records \emph{continuously and in parallel}, from the start of
the project until the final deadline, so that data accrues without interruption throughout
the study. Collecting from three participants additionally enables a cross-person evaluation
of model generalisation (Section~\ref{sec:eval}). Table~\ref{tab:data_layers} summarises the
sensors, sampling frequencies, and rationale for each layer.

\begin{table}[t]
\caption{Dual-modal data collection architecture.}\label{tab:data_layers}
\begin{tabular}{@{}p{2.6cm}p{2.8cm}p{1.8cm}p{4.6cm}@{}}
\toprule
\textbf{Data Layer} & \textbf{Primary Sensor} & \textbf{Frequency} & \textbf{Rationale} \\
\midrule
Kinematic           & iPhone Accelerometer \& Gyroscope (Sensor Logger) & 10 Hz &
    Tracks macroscopic movements and phone placement. Prolonged physical stillness indicates
    a sedentary position, raising the baseline probability of phone engagement, while sudden
    spikes capture pickup gestures. \\
\addlinespace
System \& Target    & iOS Shortcuts (event-driven) & On launch &
    Automates ground-truth timestamp logging for target app opens (TikTok, Instagram) and
    for auxiliary app opens used as behavioural features. \\
\bottomrule
\end{tabular}
\end{table}

\paragraph{Kinematic Layer.}
The iPhone's internal accelerometer and gyroscope are logged at 10\,Hz using the Sensor
Logger application, which supports persistent background recording so that motion is
captured even when the screen is locked or another app is in the foreground. The primary
derived signal is \emph{phone stillness duration}: a running counter measuring how long the
composite gyroscope variance has remained below a static threshold. Extended stillness is a
behavioural precursor to distraction-seeking. A low sampling frequency is sufficient here:
it captures pickup gestures and restlessness while keeping data volume manageable.

\paragraph{System and Target Layer.}
iOS Shortcuts automate the logging of two event types to a shared CSV file on iCloud Drive:
(1) every launch of TikTok or Instagram (the classification target), and (2) every launch of
an auxiliary (non-target) application, tagged at log time with one of four categories:
\emph{Search} (LLM assistants such as Gemini and Claude, and web browsers),
\emph{Social} (e.g.\ WhatsApp and Snapchat),
\emph{Entertainment} (e.g.\ Netflix and YouTube), and
\emph{Other} (any remaining application). These auxiliary opens serve as proxies for
attention fragmentation. The Phone app is excluded, as calling behaviour is not expected to
correlate with scrolling onset. This provides precise, automated ground-truth labels and
behavioural features without manual annotation. Because iOS automations can only fire on app
\emph{open} (not \emph{close}), all event-driven features are anchored on open events.

\subsection{Engineered Features}

Raw sensor streams are aggregated into three-minute tabular blocks to produce the following
features.

\paragraph{Kinematic features.}
\begin{itemize}
    \item \texttt{phone\_stillness\_duration} --- seconds for which the composite gyroscope
          variance stayed below a strict static threshold.
    \item \texttt{gyro\_max\_vector} --- magnitude of the single largest angular-velocity
          spike in the window, capturing rapid actions such as picking up the device.
    \item \texttt{gyro\_variance} --- variance of the gyroscope signal over the window, a
          proxy for user restlessness.
\end{itemize}

\paragraph{System and interaction features.}
\begin{itemize}
    \item \texttt{other\_apps\_opened\_count} --- total auxiliary opens in the window, plus
          per-category counts (\emph{Search}, \emph{Social}, \emph{Entertainment},
          \emph{Other}) as finer-grained attention-fragmentation signals.
    \item \texttt{app\_switch\_frequency} --- rate of transitions between distinct app
          categories within the block.
\end{itemize}

\paragraph{Contextual and temporal anchors.}
\begin{itemize}
    \item \texttt{time\_of\_day\_sin} / \texttt{time\_of\_day\_cos} --- cyclical encoding of
          the 24-hour clock, mapping circadian vulnerability without a midnight discontinuity.
    \item \texttt{time\_since\_last\_target\_open} --- seconds since the last target app-open
          event, capturing the behavioural refractory period between doomscrolling bouts
          (anchored on opens, since closes cannot be logged).
\end{itemize}

\subsection{Evaluation Strategy}\label{sec:eval}

The evaluation is designed to measure \emph{generalisation}, never random time-point sampling
within a single continuous recording: adjacent three-minute blocks are autocorrelated, so a
random split would leak information between train and test sets.

\begin{itemize}
    \item \textbf{RQ1 (Predictability).} Train and test on the same user, but split by
          \emph{recording session / day} --- earlier sessions train, held-out later sessions
          test --- preventing temporal leakage between neighbouring windows.
    \item \textbf{RQ2 (Transferability).} Use a leave-one-person-out scheme across the three
          participants: train on two users and test on the held-out third, whose data never
          appears in training. Performance is compared against that user's own within-person
          baseline to quantify the transfer gap.
    \item \textbf{Metrics.} Because Class~1 (susceptible) blocks are rare, we report PR-AUC,
          precision, recall, and F1 with a confusion matrix, relative to a majority-class
          baseline, rather than raw accuracy.
\end{itemize}

% -----------------------------------------------------------------------
\section{Data Collection and Exploratory Analysis}\label{sec:eda}
% -----------------------------------------------------------------------

% TODO (Deadline 1): Summarise collected data. Report data volumes, class
% distributions, basic statistics per feature, and time-series plots.
% Discuss resolution choices (e.g. window size, resampling rate).


% -----------------------------------------------------------------------
\section{Noise Removal and Missing Value Handling}\label{sec:noise}
% -----------------------------------------------------------------------

% TODO (Deadline 1): Document preprocessing steps. Justify chosen
% interpolation, smoothing, or imputation techniques with reference to
% Chapter 3 of the book. Discuss edge cases (sensor dropout, sync lag).


% -----------------------------------------------------------------------
\section{Feature Engineering}\label{sec:features}
% -----------------------------------------------------------------------

% TODO (Deadline 2): Expand on the feature set defined in Section 1.
% Include ablation analysis or feature importance scores to justify the
% final feature set. Reference Chapter 4 of the book.


% -----------------------------------------------------------------------
\section{Classical Machine Learning}\label{sec:classical}
% -----------------------------------------------------------------------

% TODO (Deadline 2): Define train/test split. Apply classifiers (e.g.
% XGBoost, Random Forest, SVM). Report hyperparameter optimisation (e.g.
% grid search, cross-validation). Discuss results with reference to
% Chapter 7 of the book.


% -----------------------------------------------------------------------
\section{Deep Learning}\label{sec:deep}
% -----------------------------------------------------------------------

% TODO (Final): Apply a temporal deep learning model (e.g. LSTM, TCN) on
% the raw 10 Hz sequences. Justify architecture and hyperparameter choices.
% Compare against classical results from Section~\ref{sec:classical}.


% -----------------------------------------------------------------------
\section{Conclusion and Critical Reflection}\label{sec:conclusion}
% -----------------------------------------------------------------------

% TODO (Final): Summarise findings across both modelling approaches.
% Critically reflect on limitations: data quantity, single-user bias,
% sensor reliability, label quality, and potential confounds.


% ---- Bibliography ----
\begin{thebibliography}{8}



\end{thebibliography}

\end{document}
